\subsection{Гироскопическая стабилизация.}
При наличии гироскопических сил линеаризованные уравнения Лагранжа имеют вид
\[
    A\ddot{\bt{q}} + G\dot{\bt{q}} +C\bt{q} = 0.
\]
Где матрица $G$ -- кососимметрическая (гироскопические силы имеют нулевую мощность).
\begin{definition}
    Степенью неустойчивости (по Пуанкаре) будем называть число отрицательных корней полинома $\det (C - \lambda A)$.
\end{definition}
\begin{theorem}[Кельвин, Четаев]
    Пусть есть некоторая консервативная система с нулевым положением равновесия.
    Тогда справедливы следующие утверждения:
    \begin{itemize}
        \item Если степень неустойчивости системы не кратна двум, то система не может быть стабилизирована добавлением гироскопических сил.
        \item Если степень неустойчивости системы кратна двум, то система допускает гироскопическую стабилизацию.
    \end{itemize}
\end{theorem}
\begin{proof}
    Характеристический полином системы имеет вид
    \[
        P(\lambda) = \det A \lambda^2 + B \lambda + C.
    \]
    Пусть степень неустойчивости нечётна.
    Тогда $\det C < 0$ и $P(0) = \det C < 0$.
    В тоже время, поскольку $A$ -- положительно определённая матрица, то $P(+\infty) > 0$.
    Но тогда независимо от $B$ у нас есть корень $\lambda > 0$ и система будет неустойчива по теореме об устойчивости по первому приближению при произвольной линейной стабилизации. \\
    Пусть теперь степень неустойчивости чётна.
    Перейдём к нормальным координатам $\bt{\theta}$:
    \[
        \bt{q} = U\bt{\theta} \quad I = U^{T}AU \quad \Lambda = U^T C U.
    \]
    Возьмём два неустойчивых $\lambda$.
    Без ограничения общности это будут $\lambda_1$ и $\lambda_2$.
    Для гироскопической стабилизации надо найти такое $g$ чтобы полином следующей системы стал устойчивым:
    \[
        \begin{pmatrix*} \ddot{\theta_1} \\ \ddot{\theta_2}
        \end{pmatrix*} + \begin{pmatrix*} 0 & g \\ -g & 0
        \end{pmatrix*} \begin{pmatrix*} \dot{\theta_1} \\ \dot{\theta_2}
        \end{pmatrix*} + \begin{pmatrix*} \lambda_1 & 0 \\ 0 & \lambda_2
        \end{pmatrix*} \begin{pmatrix*} \theta_1 \\ \theta_2
        \end{pmatrix*} = \bt{0}.
    \]
    Запишем его:
    \[
        \det \begin{pmatrix*} \lambda^2 + \lambda_1 & g\lambda \\ -g\lambda & \lambda^2 + \lambda_2
        \end{pmatrix*} = 0.
    \]
    После раскрытия скобочек и приведения подобных слагаемых он примет вид
    \[
        \lambda^4 + \lambda^2(\lambda_1 + \lambda_2 + g^2) + \lambda_1\lambda_2 = 0.
    \]
    Поскольку корни записываются как
    \[
        \lambda^2 = \dfrac{-(\lambda_1 + \lambda_2 + g) \pm \sqrt{D}}{2}.
    \]
    То для стабилизации полинома нам достаточно чтобы $\lambda_1 + \lambda_2 + g^2 > 0$ и $D > 0$.
    В силу неограниченности $g^2$ сверху это возможно и $\lambda^2 < 0$, а следовательно положение равновесие станет устойчивым в линейном приближении.
\end{proof}
\subsection{Вынужденные колебания: нерезонансный случай.}
Пусть у нас есть линейная система, которая движется под действием внешней гармонической силы:
\[
    A\ddot{\bt{q}} + B\dot{\bt{q}} + C\bt{q} = \bt{Q}, \quad \bt{Q} = \bt{a}\cos wt.
\]
Решением неоднородной системы является
\[
    \bt{q}(t) = \bt{q}^{\text{одн}}(t) + \bt{q}^{\text{част}}(t).
\]
Поскольку решение однородной системы мы уже искать умеем, то найдём частное решение неоднородной.
Будем его искать в виде $\bt{u}e^{iwt}$.
После подстановки в уравнение получим
\[
    (-A w^2 + i Bw + C)\bt{u}e^{iwt} = \bt{a}e^{iwt}.
\]
Обозначая матрицу $-A w^2 + i Bw + C$ как $D = W^{-1}$ мы получаем что частное решение будет равно
\[
    \bt{u} = W \bt{a}.
\]
Элементы $W$ равны
\[
    W_{jk} = \dfrac{\Delta^{jk} D}{\det D}.
\]
И тогда $j$-ая координата частного решения выражается как
\[
    q^{\text{част}}_j(t) = \sum\limits_{k}|W_{jk}(w)| \cdot a_k \cdot \cos (wt + \arg W_{jk}).
\]
Есть некоторые стандартные называния для параметров вынужденных колебаний:
\begin{itemize}
    \item $W_{jk}$ называется амплитудно-фазовой характеристикой,
    \item $|W_{jk}|$ называется амплитудно-частотной характеристикой,
    \item $\arg W_{jk}$ называется фазово-частотной характеристикой.
\end{itemize}
